% WACV 2027 Paper Template
% based on the ICCV 2025 template (https://media.eventhosts.cc/Conferences/ICCV2025/ICCV2025-Author-Kit-Feb.zip) with
% WACV-specific details (e.g., 2 tracks) from the WACV 2025 template (https://www.dropbox.com/scl/fi/su44zgdhrzik26p2xu37k/WACV-2025-Author-Kit-Template.zip?rlkey=5qcfimjhxnmx3wlyk7yhk8wg7&dl=0)

\documentclass[10pt,twocolumn,letterpaper]{article}

%%%%%%%%% PAPER TYPE  - PLEASE UPDATE FOR FINAL VERSION
\usepackage[review,algorithms]{wacv}      % To produce the REVIEW version for the algorithms track
% \usepackage[review,applications]{wacv}
% \usepackage[review,datasets]{wacv}
% \usepackage{wacv}              % To produce the CAMERA-READY version
%\usepackage[pagenumbers]{wacv} % To force page numbers, e.g. for an arXiv version

% preamble.tex — additional packages for PhysDrop paper

\usepackage{amsmath}
\usepackage{amssymb}
\usepackage{booktabs}
\usepackage{multirow}
\usepackage{graphicx}
% subcaption already loaded by wacv.sty
\usepackage{xcolor}
\usepackage{url}
\usepackage{microtype}
\usepackage{siunitx}
\usepackage{enumitem}
\usepackage{tablefootnote}

% Convenient math macros
\newcommand{\betamax}{\ensuremath{\beta_{\max}}}
\newcommand{\cor}{\ensuremath{\eta}}
\newcommand{\Dzero}{\ensuremath{D_0}}
\newcommand{\Uzero}{\ensuremath{U_0}}
\newcommand{\We}{\ensuremath{\mathrm{We}}}
\newcommand{\Oh}{\ensuremath{\mathrm{Oh}}}

% Table helpers
\newcommand{\best}[1]{\textbf{#1}}
\newcommand{\second}[1]{\underline{#1}}
\definecolor{failred}{rgb}{0.8,0.1,0.1}
\newcommand{\fail}[1]{{\color{failred}#1}}


\definecolor{wacvblue}{rgb}{0.21,0.49,0.74}
\usepackage[pagebackref,breaklinks,colorlinks,allcolors=wacvblue]{hyperref}

%%%%%%%%% PAPER ID  - PLEASE UPDATE
\def\wacvPaperID{*****}
\def\confName{WACV}
\def\confYear{2027}

%%%%%%%%% TITLE
\title{PhysDrop: Benchmarking Vision-Language Models on\\
Physical Reasoning from Scientific High-Speed Droplet Impact Video}

%%%%%%%%% AUTHORS
\author{Rasul Khanbayov\\
% Institution, Department\\
% City, Country\\
{\tt\small barhdadirayan01@gmail.com}
}

\begin{document}
\maketitle

% ─────────────────────────────────────────────────────────────────────────────
\begin{abstract}
% ─────────────────────────────────────────────────────────────────────────────
Scientific video analysis in experimental fluid mechanics relies on
hand-tuned classical computer-vision pipelines that require domain expertise
and manual re-calibration for each experimental setup.
We ask whether modern vision-language models (VLMs) can replace or
extend these pipelines using a novel dataset of \textbf{99 unpublished
high-speed shadowgraphy videos} of $\sim$\SI{2}{mm} droplets impacting a
superhydrophobic glass surface across 8 experimental conditions (pure water,
three surfactant types above and below CMC, and cellulose-acetate
nanoparticles with and without surfactant).

We introduce \textbf{PhysDrop}, a benchmark of 1{,}211 annotated entries
covering two complementary task families: (1)~\emph{measurement tasks}—impact
phase classification and quantitative localisation (centroid, radius, spread
width); and (2)~\emph{physical reasoning}—fluid composition classification
directly from observed visual impact dynamics, requiring no labels beyond the
experimental design.

We evaluate five frontier VLMs zero-shot (Gemini 2.0 Flash, GPT-4o, GPT-4o-mini,
Claude Sonnet 4.5) and one domain-fine-tuned model (Qwen2.5-VL-7B LoRA,
trained on 695 labelled frames).
On measurement tasks, zero-shot phase accuracy plateaus at 46–58\%;
fine-tuning reduces radius MAE to \SI{0.12}{mm} but does not generalise to
out-of-distribution visual conditions.
On fluid composition classification,
\textit{every} frontier VLM scores \textit{below} the 25\% chance baseline:
the best model (Gemini) achieves 21.3\%, while all GPT models collapse to
predicting a single class.
In contrast, a simple $k$-NN classifier trained on two classical-CV features
(contact time and spread width) achieves \textbf{66.0\%}—a \textbf{44.7
percentage-point} gap over the best VLM.
Our results show that current VLMs learn visual patterns but cannot perform
the physical reasoning needed to identify hidden material properties from
observed dynamics.
\end{abstract}

% ─────────────────────────────────────────────────────────────────────────────
\section{Introduction}
\label{sec:intro}
% ─────────────────────────────────────────────────────────────────────────────

High-speed video is a primary measurement instrument in experimental fluid
mechanics. A researcher records a droplet impacting a surface at several
thousand frames per second, then runs a computer-vision pipeline to extract
quantitative observables: impact velocity $\Uzero$, pre-impact diameter
$\Dzero$, maximum spreading factor $\betamax = D_{\max}/\Dzero$, and
coefficient of restitution $\cor = U_{\text{rebound}}/\Uzero$.
These quantities govern wetting, coating, and spray
processes~\cite{yarin2006drop,richard2002surface}.

Current pipelines are hand-tuned classical CV: binarisation thresholds,
HoughCircles parameters, and contact-line detectors must be adjusted
per-video-setup~\cite{abbot2025nanoparticles,jereb2025droplet}.
More capable alternatives—segment-anything models~\cite{sam2_2024} and
vision-language models (VLMs)~\cite{qwen25vl,gpt4o,gemini2flash}—have
recently shown strong performance on diverse visual tasks.
This raises a natural question: \emph{can VLMs replace or extend classical
pipelines for scientific video measurement?}

Beyond measurement, high-speed video contains richer information.
The visual dynamics of impact—how far a droplet spreads, how quickly it
retracts, whether it bounces or deposits—encode the fluid's
surface tension, viscosity, and particle content.
A skilled experimentalist can identify from a single video whether a droplet
contains surfactant or nanoparticles.
Can a VLM do the same?

We build \textbf{PhysDrop}, the first benchmark for VLMs on scientific
physical video reasoning, comprising 99 high-speed droplet videos across
8 experimental conditions and two task families:
\begin{itemize}[leftmargin=*, topsep=2pt, itemsep=1pt]
  \item \textbf{Task 1/4 — Measurement:} per-frame phase classification
    (falling / spreading / rebounding) and quantitative localisation
    (centroid, radius, spread width).
  \item \textbf{Task 6 — Physical Reasoning:} given a 3-frame summary
    (pre-impact, max-spread, post-rebound), classify the fluid composition
    from four classes: pure water, surfactant only,
    CA nanoparticles + surfactant, washed CA nanoparticles.
    Ground truth labels come directly from the experimental design—no extra
    annotation is needed.
\end{itemize}

\noindent\textbf{Key findings:}
\begin{enumerate}[leftmargin=*, topsep=2pt, itemsep=1pt]
  \item Zero-shot VLMs achieve 46–58\% on 3-class phase detection (chance
    33\%), with no model exceeding 58\%.
  \item Fine-tuning Qwen2.5-VL-7B with 695 frames reduces radius MAE from
    0.23–0.38~mm to 0.12~mm (best across all models), but phase accuracy
    drops to 50\% on out-of-distribution videos.
  \item \emph{All} VLMs—including the fine-tuned model—score below 25\% chance
    on fluid composition classification. Both GPT models predict a single class
    for every video; Gemini predicts a different single class.
  \item A $k$-NN classifier on two automatically-extracted CV features achieves
    66.0\%—a 44.7 pp gap—demonstrating that physical features, not visual
    patterns, are key for fluid characterisation.
\end{enumerate}

These results establish that VLMs can learn to \emph{measure} physical scenes
with fine-tuning, but cannot \emph{reason} about the underlying physics from
visual appearance alone.
This gap defines the open problem for AI in experimental science.

% ─────────────────────────────────────────────────────────────────────────────
\section{Related Work}
\label{sec:related}
% ─────────────────────────────────────────────────────────────────────────────

\paragraph{Classical CV for droplet analysis.}
Abbot~\etal~\cite{abbot2025nanoparticles} analyse silica nanoparticle droplets
with a custom MATLAB pipeline (binarisation $\to$ boundary detection $\to$
circle fitting), concluding that additive-free NPs do not affect impact dynamics.
Jereb~\etal~\cite{jereb2025droplet} train Gaussian Process Regression models
on scalars extracted by classical contour detection ($\betamax$: $R^2{=}0.97$),
but do not address image-level AI. Neither work asks whether foundation models
can automate the extraction pipeline or infer fluid properties directly.

\paragraph{VLM benchmarks.}
Existing VLM evaluation suites target static images with textbook
knowledge~\cite{mmmu2024,scienceqa2022,mathvista2024} or scientific
charts~\cite{charxiv2024,sciverse2024}. No existing benchmark evaluates
VLMs on dynamic physical processes captured in high-speed video, nor tests
whether models can infer hidden physical properties from observed dynamics.

\paragraph{AI for Science.}
Foundation models have transformed structural biology~\cite{jumper2021alphafold}
and materials discovery~\cite{ai4science2023}. Video
VLMs~\cite{videollava2024,internvl2024} extend multimodal reasoning to dynamic
scenes, but have not been evaluated on quantitative scientific measurement or
physical property inference from video.

\paragraph{Domain fine-tuning.}
LoRA~\cite{lora2022} enables efficient adaptation of large VLMs to narrow tasks
with hundreds of labelled examples. We show that fine-tuning recovers
near-CV accuracy on measurement tasks within the training distribution, but
fails to generalise to new visual conditions or to out-of-training-set tasks.

% ─────────────────────────────────────────────────────────────────────────────
\section{Dataset}
\label{sec:dataset}
% ─────────────────────────────────────────────────────────────────────────────

\subsection{Experimental setup}

Droplets of $4\,\mu$L were dispensed from a fixed height onto a
superhydrophobic borosilicate glass surface and imaged with a
Photron FASTCAM Nova S6 at 2{,}997\,fps (1280$\times$512\,px).
Two impact velocities were used: $U_0 \approx \SI{1.14}{m/s}$ (drop height
6.5\,cm) and $U_0 \approx \SI{1.63}{m/s}$ (13.5\,cm), corresponding to Weber
numbers $\We = \rho U_0^2 \Dzero / \sigma \approx 36$ and $73$, respectively.
Scale calibration gives 65.6--66.0\,px/mm.

\subsection{Experimental conditions and fluid classes}

Eight sample types were recorded across five recording sessions, spanning two
impact velocities (Table~\ref{tab:dataset}).
Samples divide naturally into four fluid classes (Task~6 labels):

\begin{description}[leftmargin=0.5cm, topsep=2pt, itemsep=1pt]
  \item[\textbf{A — Pure water}] DI water with or without inert fluorescent
    tracer (Nile Red). No additives that affect surface tension.
  \item[\textbf{B — Surfactant only}] Sodium dodecyl sulfate (SDS),
    Triton X-100 (TX-100), or cocoglycoside (CG) at concentrations above
    and below the critical micelle concentration (CMC).
  \item[\textbf{C — CA + surfactant}] Cellulose acetate (CA) nanoparticles
    synthesised \emph{in the presence of} surfactant; free surfactant
    remains in the droplet.
  \item[\textbf{D — Washed CA}] CA nanoparticles prepared with surfactant
    then washed; no free surfactant. Particle surface chemistry is
    retained but bulk surfactant is removed.
\end{description}

\begin{table}[t]
\centering
\caption{PhysDrop dataset summary. $N_v$: video count; $N_{t6}$: Task 6
benchmark entries (one per video).}
\label{tab:dataset}
\small
\begin{tabular}{lccc}
\toprule
Session & Conditions & $N_v$ & $N_{t6}$ \\
\midrule
02182026 & A, B (3 surf.), C (3 surf.) & 32 & 26 \\
03242026 & B (6 conc.), D & 23 & 23 \\
05052026 & B, C (CG, TX, SDS above CMC) & 19 & 19 \\
05112026 & A, D (high vel.) & 7 & 7 \\
05122026 & B, C (3 surf., high vel.) & 18 & 18 \\
\midrule
\textbf{Total} & \textbf{8 conditions} & \textbf{99} & \textbf{93} \\
\bottomrule
\end{tabular}
\end{table}

All videos are unpublished and not included in any public pre-training corpus.

\subsection{Classical CV pipeline}

We extract ground-truth physical parameters using a three-stage pipeline
implemented in Python/OpenCV~\cite{opencv}:
\textbf{(1) Impact detection:} the frame of maximum pixel-difference in the
lower half of the frame;
\textbf{(2) Droplet localisation:} HoughCircles~\cite{opencv} on pre-impact
and rebound frames to obtain centroid and radius;
\textbf{(3) Spread measurement:} background-subtracted contact width at the
surface row during the spreading phase.
This yields: $\Dzero$, $\Uzero$, $\betamax$, contact time $\tau$, and—for
a subset—$\cor$.
The same pipeline generates ground-truth labels for the measurement benchmark
(Tasks~1/4) and physical features for the classical-CV Task~6 classifier.

% ─────────────────────────────────────────────────────────────────────────────
\section{PhysDrop Benchmark}
\label{sec:benchmark}
% ─────────────────────────────────────────────────────────────────────────────

\subsection{Task 1/4 — Phase detection and measurement}
\label{sec:task14}

For each video, we extract 12 frames: 5 pre-impact (\emph{falling}), 4
during contact (\emph{spreading}), and 3 post-liftoff (\emph{rebounding}).
Each frame is paired with a prompt asking the VLM to:
(i)~classify the impact phase (falling / spreading / rebounding);
(ii)~if falling or rebounding, estimate the droplet centroid $(c_x, c_y)$
and radius $r$ in pixels;
(iii)~if spreading, estimate the horizontal contact width.
Ground-truth is provided by the classical CV pipeline.
The benchmark contains 1{,}117 single-frame entries.

\paragraph{Metrics.}
Phase accuracy (\%), centroid MAE $(c_x, c_y)$, radius MAE, and spread-width
MAE, all reported in mm using the per-video calibration.
Physics accuracy: VLM-derived $\betamax = w_{\text{spread}} / (2r)$ versus
CV-derived $\betamax$.

\subsection{Task 6 — Fluid composition classification}
\label{sec:task6}

Each video generates \emph{one} Task 6 entry consisting of three frames:
(1)~5 frames before impact; (2)~10 frames after impact (near maximum spread);
(3)~8 frames after liftoff (post-rebound).
The prompt asks the VLM to classify the fluid composition as A/B/C/D
based solely on the visual impact dynamics, without any numerical
measurements in the prompt.
Ground-truth labels come from the experimental design; no manual annotation
is required.
The benchmark contains 94 Task~6 entries (class distribution:
$n_A{=}11$, $n_B{=}24$, $n_C{=}35$, $n_D{=}24$).
Chance accuracy is 25\% (4-class).

\paragraph{Classical CV baseline.}
For videos with complete CV features ($n{=}50$), we train a $k$-nearest-
neighbour classifier ($k{=}3$, Euclidean distance in standardised feature
space) using \emph{leave-one-out cross-validation} on two features:
contact time $\tau$ (ms) and maximum spread width (mm).
These features are automatically extracted by the classical pipeline and
require no additional labelling.

% ─────────────────────────────────────────────────────────────────────────────
\section{Experiments}
\label{sec:experiments}
% ─────────────────────────────────────────────────────────────────────────────

\subsection{Models evaluated}

\paragraph{Zero-shot VLMs (via OpenRouter API).}
We evaluate four frontier VLMs in zero-shot mode with greedy decoding and
a structured JSON output prompt:
Gemini 2.0 Flash~\cite{gemini2flash},
GPT-4o and GPT-4o-mini~\cite{gpt4o},
and Claude Sonnet 4.5~\cite{claude35sonnet}.

\paragraph{Fine-tuned model.}
We fine-tune Qwen2.5-VL-7B-Instruct~\cite{qwen25vl} with LoRA
($r{=}16$, $\alpha{=}32$)~\cite{lora2022,peft} on 695 frames from the
02182026 recording session, labelled with phase, centroid, radius, and spread
width. Training uses the Task~1/4 prompt only. The adapter is loaded on an
NVIDIA A100 80\,GB GPU.

\subsection{Task 1/4 results}
\label{sec:results14}

Table~\ref{tab:task14} summarises phase accuracy and measurement error for all
models on the 1{,}117-entry benchmark.

\begin{table*}[t]
\centering
\caption{Task 1/4 results — phase detection and quantitative measurement.
All MAE values in mm using per-video calibration (65.6–66.0\,px/mm).
Best in \textbf{bold}, second \underline{underlined}. $^\dagger$Fine-tuned model trained
only on 02182026 session; spread\_MAE inflated by frame-width fallback on
out-of-distribution spreading frames (see text).}
\label{tab:task14}
\small
\begin{tabular}{lcccccc}
\toprule
Model & Phase Acc.\,(\%) & $c_x$ MAE & $c_y$ MAE & Radius MAE & Spread MAE & Null\,(\%) \\
\midrule
Chance (3-class)            & 33.3 & — & — & — & — & — \\
\midrule
GPT-4o-mini~\cite{gpt4o}    & 46.6 & 2.20 & 1.19 & 0.38 & 4.24 & 0.0 \\
GPT-4o~\cite{gpt4o}         & 50.9 & 1.83 & 1.11 & 0.30 & 2.22 & 0.0 \\
Fine-tuned Qwen LoRA$^\dagger$ & 50.0 & \best{0.37} & \best{0.26} & \best{0.12} & 15.17$^\dagger$ & 2.4 \\
Gemini 2.0 Flash~\cite{gemini2flash} & \best{57.6} & 1.54 & 0.80 & \second{0.23} & 1.38 & 0.1 \\
Claude Sonnet 4.5~\cite{claude35sonnet} & \best{57.6} & 0.70 & \second{0.60} & 0.31 & \best{0.80} & 0.0 \\
\bottomrule
\end{tabular}
\end{table*}

\paragraph{Phase detection.}
No model exceeds 58\% on 3-class phase detection (chance: 33\%).
Gemini and Claude Sonnet 4.5 are tied at 57.6\%.
GPT-4o improves over GPT-4o-mini by only 4 pp despite substantially higher
compute cost, suggesting diminishing returns of model scale on this task.

\paragraph{Spatial localisation.}
The fine-tuned Qwen model achieves the best centroid ($c_x$ MAE: 0.37\,mm)
and radius MAE (0.12\,mm) across all models, demonstrating that fine-tuning
effectively transfers localisation accuracy.
Claude Sonnet 4.5 achieves the best spread-width MAE (0.80\,mm) among
zero-shot models, considerably better than GPT-4o (2.22\,mm) and
GPT-4o-mini (4.24\,mm).

\paragraph{Fine-tuned generalisation gap.}
The Qwen LoRA model achieves 93.3\% phase accuracy on its in-distribution
validation set (78 frames from 02182026) but drops to 50.0\% on the full
benchmark.
Per-session analysis reveals the collapse is driven by the new-experiments
sessions (05112026: 28.7\%; 05122026: 37.7\%), where the Nile Red fluorescent
tracer creates a distinctive ring-shaped droplet appearance unseen during
training.
This demonstrates a fundamental generalisation limitation of domain fine-tuning
under visual distribution shift.

The spread-width MAE of 15.17\,mm is caused by the model outputting the frame
width (1{,}279\,px) as a fallback for spreading frames it cannot segment—a
degenerate but deterministic failure mode. For in-distribution spreading frames
in 02182026, spread-width predictions are accurate.

\subsection{Task 6 results}
\label{sec:results6}

Table~\ref{tab:task6} shows fluid composition classification accuracy for all
models. Figure~\ref{fig:confusion} visualises the per-class prediction
distribution for each model.

\begin{table}[t]
\centering
\caption{Task 6 — fluid composition classification accuracy (\%).
4-class problem; chance = 25\%.
Per-class: A = pure water, B = surfactant only, C = CA+surfactant, D = washed CA.
\fail{Red}: below chance baseline.}
\label{tab:task6}
\small
\setlength{\tabcolsep}{3pt}
\begin{tabular}{lcccccl}
\toprule
Model & Overall & A & B & C & D & Bias \\
\midrule
Chance baseline      & 25.0 & 25 & 25 & 25 & 25 & — \\
\midrule
GPT-4o               & \fail{10.6} & 90.9 & 0 & 0 & 0 & Always A \\
GPT-4o-mini          & \fail{11.7} & 100 & 0 & 0 & 0 & Always A \\
FT-Qwen              & \fail{11.7} & 100 & 0 & 0 & 0 & Always A \\
Claude 4.5           & \fail{18.1} & 45.5 & 8.3 & 2.9 & 37.5 & Mixed A/D \\
Gemini 2.0F          & \fail{21.3} & 18.2 & 0 & 0 & 75.0 & Always D \\
\midrule
\textbf{CV (kNN-3)}  & \best{66.0} & 100 & 80.0 & 66.7 & 52.4 & Balanced \\
\bottomrule
\end{tabular}
\end{table}

\begin{figure*}[t]
\centering
\includegraphics[width=\linewidth]{fig_confusion}
\caption{Confusion matrices for Task~6 fluid composition classification
(row = true label, column = predicted label; cell values are counts; shading
is row-normalised).
A = pure water, B = surfactant only, C = CA+surfactant, D = washed CA.
GPT-4o, GPT-4o-mini, and Fine-tuned Qwen collapse entirely to class A;
Gemini collapses to class D; Claude Sonnet 4.5 is more distributed but still
below the 25\% chance baseline.
\textit{No VLM correctly identifies any class~B or class~C video.}}
\label{fig:confusion}
\end{figure*}

\paragraph{Universal VLM failure.}
Every VLM—zero-shot and fine-tuned—scores below the 25\% chance baseline.
The failure is not random: models develop strong class biases.
GPT-4o and GPT-4o-mini predict class A (pure water) for \emph{all 94} videos;
Gemini predicts class D (washed CA) for 77/94 videos;
Claude is more distributed but still below chance.
Notably, scaling GPT-4o-mini to GPT-4o \emph{decreases} Task~6 accuracy
(11.7\% $\to$ 10.6\%), confirming that model capability does not translate
to physical reasoning.
The fine-tuned model, trained exclusively on Task~1/4 prompts, inherits the
GPT-style A-bias upon encountering the Task~6 prompt, scoring identically to
GPT-4o-mini (11.7\%).

\paragraph{Classical CV advantage.}
The $k$-NN classifier using only $\tau$ and $w_{\text{spread}}$ reaches
66.0\% LOO accuracy—a \textbf{44.7 pp gap} over the best VLM (Gemini, 21.3\%).
Per-class accuracy is 100\% on pure water, 80\% on surfactant-only, 66.7\% on
CA+surfactant, and 52.4\% on washed CA.
The harder classes (C and D) differ physically by whether free surfactant is
present in the droplet bulk—a subtle distinction that classical features
encode implicitly through the contact-time and spread dynamics.

\subsection{Why VLMs fail at Task 6}

\paragraph{Single-class collapse.}
The consistent collapse to a single class (A or D depending on model)
suggests VLMs are applying a \emph{prior} over class labels rather than
extracting task-relevant visual cues from the impact sequence.
Absent clear visual discriminators in their training data for this domain,
models default to the most frequent or most ``natural'' label.

\paragraph{No temporal integration.}
The three-frame Task~6 prompt spans $\sim$50 frames of dynamics.
VLMs likely process each frame independently and average the evidence
rather than tracking continuous physical quantities (spread rate, retraction
speed, satellite formation) that encode fluid properties.

\paragraph{Classical CV encodes physics.}
Contact time $\tau$ is related to the capillary time
$\tau_c = \sqrt{\rho \Dzero^3 / \sigma}$ (surface tension $\sigma$);
spread width encodes $\betamax \propto \We^{1/4}$ for low-viscosity
fluids~\cite{yarin2006drop}.
Surfactants reduce $\sigma$, systematically increasing $\tau$ and
$\betamax$ for class B/C, while washed NPs (class D) behave more like
water (class A) in bulk but retain modified wetting kinetics.
The $k$-NN classifier discovers these class-separating directions in
$(\tau, w_{\text{spread}})$ space; VLMs, processing raw pixels without
physics priors, cannot.

% ─────────────────────────────────────────────────────────────────────────────
\section{Discussion}
\label{sec:discussion}
% ─────────────────────────────────────────────────────────────────────────────

\paragraph{Measurement vs.\ physical reasoning.}
Our benchmark reveals a sharp bifurcation in VLM capability:
with fine-tuning, a 7B model can \emph{locate} a $\sim$2\,mm droplet
with sub-pixel accuracy and classify its impact phase with 93\% accuracy
(in-distribution).
Yet the same model—and all frontier zero-shot VLMs—completely fail to infer
what is \emph{inside} the droplet from how it moves.
This separation between perception and physics-grounded reasoning defines the
current frontier of AI in experimental science.

\paragraph{Implications for lab automation.}
Classical CV pipelines require per-setup calibration but encode
domain physics through their feature choices.
VLMs promise automation without calibration, but our results show they do
not yet provide a drop-in replacement, particularly for inference tasks
that require understanding physical causality.
A promising direction is hybrid pipelines that use VLMs for robust
localisation and segmentation, while delegating physical interpretation
to physics-aware models or classical extractors.

\paragraph{Dataset velocity and generalisation.}
The two-velocity design (1.14 and 1.63\,m/s) reveals that fine-tuned
model accuracy degrades severely under visual distribution shift
introduced by the fluorescent tracer (Nile Red), even though the physical
task is identical.
Future work should investigate domain-adaptive approaches and the role of
data diversity in scientific video fine-tuning.

\paragraph{Limitations.}
The classical CV classifier is evaluated on $n{=}50$ videos from two sessions
only; the VLM evaluation covers all 94 Task~6 entries.
This discrepancy should be addressed by extending the classical pipeline to
new-experiments sessions.
COR is available for only 10 videos, precluding its use as a classifier
feature despite its strong physical relevance.
Adding COR would likely further improve the classical baseline.

% ─────────────────────────────────────────────────────────────────────────────
\section{Conclusion}
\label{sec:conclusion}
% ─────────────────────────────────────────────────────────────────────────────

We introduced PhysDrop, a benchmark for evaluating vision-language models on
scientific high-speed video of droplet impact, covering measurement tasks
(phase, centroid, radius, spread) and a physical reasoning task (fluid
composition classification).
Evaluating five VLMs across 1{,}211 benchmark entries, we find:
fine-tuning with $<$700 frames substantially improves measurement accuracy
but does not generalise to novel visual conditions;
and all VLMs—including fine-tuned—score below the 25\% chance baseline on
fluid classification, while a 2-feature $k$-NN achieves 66\%.

The 44.7\,pp gap between classical CV and the best VLM on physical reasoning
is our central finding.
It quantifies precisely where the current generation of VLMs fails: not at
detection or localisation, but at understanding the physical meaning of what
they see.
Closing this gap—whether through physics-constrained decoding, synthetic
training data from simulators, or novel architectures—is the challenge for
the next generation of AI for experimental science.

\paragraph{Acknowledgements.}
% TODO: Add acknowledgements (supervisor, funding, lab).

{\small
\bibliographystyle{ieeenat_fullname}
\bibliography{main}
}

\end{document}
